\subsection{Interaction structure at fixed input-sensitivity rank}
\label{note:interaction_structure}

The preceding linear experiment allowed a dense trainable input map, which
can compensate for generic route rotations. Let $\Phi$ contain $K$
orthonormal route columns, write $W=\Phi V$ with
$V\in\mathbb R^{K\times8}$, and let $Q_r$ contain the $r$ active
context columns of $Q$. The context-visible transfer is then
$Q_r^{\mathsf T}H_T\Phi$. When this matrix
has row rank $r$, a suitable $V$ reproduces every noiseless context. Finite-time
conditioning and cost can still differ, but detailed shape then need not
restrict representation. We therefore conducted a separate mechanism study
with explicit local input access. This analysis followed the original
prospective outcomes and is an exhaustive population diagnostic, not another
prospective test of the original selector.

To impose local input access, we placed eight independent Rademacher inputs
(each equally likely to be $-1$ or $+1$) at named leaves of a binary tree.
For internal node $v$, let $u_{\mathrm L}$ and $u_{\mathrm R}$ be its
child outputs and let $a_v,b_v,c_v,d_v$ be its four trainable coefficients.
Each of the seven internal nodes computes
\begin{equation}
 u_v=a_v+b_vu_{\mathrm L}+c_vu_{\mathrm R}
       +d_vu_{\mathrm L}u_{\mathrm R}.
 \label{eq:s_multiaffine_node}
\end{equation}
Only the root is read out. Every candidate has 28 free coefficients and 14
edges, with no dense input mixing, intermediate readout or access to inputs
outside a node's subtree.
The fixed pool combines three shapes---a balanced tree, a comb and a tree
with a $3+5$ root split---with the same four predetermined leaf permutations,
giving twelve candidates. Equal edge
and coefficient counts do not imply equal cable length, latency or biological
energy cost. These signed algebraic nodes are not conductance compartments.

The primary target census contains all 105 perfect matchings $M$ and all
35 unordered four-plus-four partitions $A$ of the eight inputs:
\begin{equation}
 f_M(\bm x)=\tfrac12\sum_{(i,j)\in M}x_ix_j,
 \qquad
 f_A(\bm x)=\tfrac12\left(\prod_{i\in A}x_i+
                              \prod_{i\notin A}x_i\right).
 \label{eq:s_interaction_targets}
\end{equation}
The targets are enumerated independently of the candidate trees. Evaluating
all $2^8=256$ equally weighted inputs gives target variances 1 and $1/2$,
respectively. For the continuous derivative of each target's unique
multilinear extension,
\begin{equation}
 \E[\nabla_{\bm x}f\,\nabla_{\bm x}f^{\mathsf T}]=I_8/4.
 \label{eq:s_equal_input_sensitivity}
\end{equation}
Thus both algebraic and participation-ratio effective rank are eight, and
the complete input-sensitivity spectrum agrees. This is a property of the
target, not equality of changing model parameter-gradient or synaptic-credit
moments. Normalized MSE (NMSE) divides population squared error by the stated
target variance.

Interaction-tensor descriptions have previously been used to analyze
structured neural function classes \citep{cohen2016tensor}. The centered
constant-feature criterion below is derived for the specific local model in
Eq.~\ref{eq:s_multiaffine_node}.

\subsection{A necessary and sufficient centered-cut condition}
\label{note:centered_cut_theorem}

For each subset $U$ of input indices, define the Walsh function
$\chi_U(\bm x)=\prod_{i\in U}x_i$, with $\chi_\emptyset=1$. These
functions form an orthonormal basis under the uniform input distribution.
Writing $c_U$ for the coefficient of $\chi_U$ gives
$f=\sum_U c_U\chi_U$. If $\bm c_f$ and $\bm c_g$ collect the
coefficients of two functions $f$ and $g$, Parseval's identity gives
$\E[(f-g)^2]=\|\bm c_f-\bm c_g\|^2$. For the leaf set $S$ of an actual
subtree, reshape the target coefficients into $C_S$, with row index a subset
of $S$ and column index a subset of its complement $\bar S$. Let $C_S^\circ$ denote
this matrix after removing the row for the empty subset of $S$.

\begin{proposition}[Exact representation by a scalar multi-affine tree]
\label{prop:s_centered_tree_iff}
For a fixed binary tree with one named Rademacher input at each leaf, a target
is represented exactly by Eq.~\ref{eq:s_multiaffine_node} and a root-only
output if and only if $\operatorname{rank}(C_S^\circ)\leq1$ at every proper
subtree cut.
\end{proposition}

\begin{proof}
Every ancestor is affine in a particular child's output, and sibling inputs
are disjoint. Consequently a represented function can be written
$g=A(\bm x_{\bar S})+h_S(\bm x_S)B(\bm x_{\bar S})$. Its full cut matrix
has rank at most two; removing the empty-$S$ row leaves rank at most one.
This proves necessity.

For sufficiency, call a cut active when $C_S^\circ\ne0$. At each active
cut, choose a centered function $h_S$ whose
Fourier coefficients span the column space of $C_S^\circ$, and define
$V_S=\operatorname{span}\{1,h_S\}$. If $C_S^\circ=0$, take $h_S=0$ and
$V_S=\operatorname{span}\{1\}$. At a leaf use
$V_i=\operatorname{span}\{1,x_i\}$. The condition implies
$f\in V_S\otimes\mathcal H_{\bar S}$, where $\mathcal H$ is the complete
finite-dimensional function space on the indicated inputs.

For sibling leaf sets $L$ and $R$, the orthogonal projections onto $V_L$
and $V_R$ commute because
they act on disjoint input factors. Hence
$f\in V_L\otimes V_R\otimes\mathcal H_{\overline{L\cup R}}$.
For a non-root parent $P=L\cup R$ with nonzero centered component, write
$f=1_P\otimes A+h_P\otimes B$, with $B\ne0$. Contracting the outside
factor against a linear functional that maps $B$ to one gives
$h_P+c1_P\in V_L\otimes V_R$, where $c$ is a scalar contributed by $A$. Since the constant function also belongs to
this product space, so does $h_P$. It is therefore a linear combination of
$1,h_L,h_R,h_Lh_R$, exactly the four allowed features. Inactive parent
features are zero and cause no exception. At the root, the complete target
itself belongs to the product of its children's spaces. Induction supplies
all local coefficients and proves sufficiency.
\end{proof}

Retaining the constant feature in every subtree makes this criterion
stronger than an uncentered rank-two constraint. Let
$\mathcal T_{\rm proper}$ be the leaf sets of the tree's proper subtrees,
and let $\sigma_j(C_S^\circ)$ be the singular values in decreasing
order. The optimal low-rank approximation identity and Parseval then give
\begin{equation}
 \operatorname{MSE}(f,g)\geq
 \max_{S\in\mathcal T_{\rm proper}}
       \sum_{j\geq2}\sigma_j(C_S^\circ)^2.
 \label{eq:s_centered_cut_bound}
\end{equation}
Taking the maximum preserves a lower bound. Summing tails from overlapping
cuts instead gives a heuristic score, without the same guarantee. A positive maximum-cut bound need not equal the true optimum.
The condition applies to actual subtree cuts; the complement of a subtree
need not itself be a subtree. For the matching $(12)(34)(56)(78)$, an
even--odd root split has centered tail $0.75$, whereas the full rank-two tail
is $0.5$. For the two-quartet target with supports $1234$ and $5678$, a
five-leaf subtree $12345$ has centered tail $0.25$ despite a zero full
rank-two tail. Supplementary Fig.~\ref{fig:supp_interaction_structure}
compares both bounds with fitted population error.

The proof requires ancestors to be affine in each child's output. A general
nonlinear scalar ancestor can generate higher cut rank from one scalar
feature. Neither the necessity condition nor Eq.~\ref{eq:s_centered_cut_bound}
therefore applies automatically to a tanh tree or a conductance model.

\subsection{Population fitting, complete-tree construction and depth}
\label{note:constructive_morphology}

All twelve candidates were fitted to every primary target by 32 alternating
least-squares (ALS) sweeps with four fixed random restarts. Conditional on the
other nodes, the design for one node is
$q(\bm x)[1,u_L,u_R,u_Lu_R]$, where $q$ is the derivative of the root with
respect to that node's output. The offset is the current root output minus
$q u_v$. Sweeps alternate postorder and preorder. Hidden mean/scale gauges
are normalized with exact parent compensation. Every restart and checkpoints
at 0, 2, 8 and 32 sweeps are retained: 6,720 fits and 26,880 checkpoint rows.
The endpoint is the smallest achieved NMSE among the four restarts. The difference between this endpoint and
Eq.~\ref{eq:s_centered_cut_bound} can reflect both incomplete optimization
and a loose bound, so it is not a certified optimization gap. Full-domain fitting
has no held-out generalization interpretation.

For score-based selection among the twelve fixed candidates, values within
$10^{-10}$ are treated as ties and resolved by lexicographic candidate
identifier. A separately retained summary applies this tolerance to the raw
floating-point minima without changing any fit.
The two-sweep pilot fits every candidate; the best-fixed baseline chooses a
single candidate per family in hindsight. The uniform-random baseline is its
exact candidate expectation. Family averages and paired shape/assignment
counts are descriptive censuses, with error comparisons at tolerance
$10^{-7}$ and no sampling confidence intervals over the enumerated tasks.
The lower bound is attained within this tolerance in 1,407/1,680 endpoints.
Although some restarts develop large cancelling coefficients, every case
selected by a deterministic policy has an existing equivalent restart within
$10^{-7}$ NMSE whose coefficient norm is at most 2.645752. All original fits,
including the remaining poorly conditioned case, remain included.

To search beyond the fixed pool, we used an exploratory dynamic program
over all input subsets. Let $c(S)$ be the normalized cut tail from
Eq.~\ref{eq:s_centered_cut_bound}, with zero cost at the root and leaves.
Let $b(S)$ be the smallest achievable maximum cut cost for a subtree on
leaf set $S$, with base case $b(\{i\})=0$. Then
\begin{equation}
 b(S)=\max\left\{c(S),\min_{L\dot\cup R=S}
                            \max[b(L),b(R)]\right\}.
 \label{eq:s_tree_minimax_dp}
\end{equation}
Every binary tree decomposes into one root split and two independent
children, so this recurrence covers all trees. A second pass restricts every
cut to the globally minimal threshold, then minimizes total tail, depth and
deterministic split ties. Costs are rounded to 12 decimals for numerical
ties. This corrected two-pass version is retained separately from the first
exploratory construction and agrees with exhaustive enumeration of all 945
six-input trees on four random cost tables. At a zero threshold, the theorem
makes the minimum-depth solution the shallowest exact tree in this class;
for a positive threshold only the cut score, not fitted error, is optimized.

Dominant left singular vectors of the centered cut matrices define local
scalar features, and four-feature least squares recovers each node's
coefficients. All 105 matching and 35 quartic targets reconstruct at minimum
depth three. A separate positive control uses 24 unique seeded permutations
$\pi$ of the nested target
\begin{equation}
 f_\pi(\bm x)=\tfrac12\sum_{k\in\{2,4,6,8\}}
                         \prod_{j=1}^k x_{\pi(j)}.
\end{equation}
Its rank is eight but its spectrum
$(1/4,1/4,1/2,1/2,3/4,3/4,1,1)$ differs from the primary families; it is
constant only within this control family. Every nested target requires
depth four. An independent depth-constrained recurrence gives minimum
centered-cut NMSE bound 0.146446609407 at depth at most three and zero at
depth four. All 164 constructions have NMSE below $7\times10^{-30}$ and
maximum absolute coefficient one up to rounding
(Supplementary Fig.~\ref{fig:supp_constructive_morphology}). The algorithm
uses the complete target Fourier tensor: these are exact-target capacity
certificates, not predictions from limited calibration data.

\subsection{Finite, noisy calibration of interaction-based tree selection}
\label{note:finite_calibration_morphology}

The exact-target construction leaves open whether a useful tree can be chosen from a limited sample. We therefore tested selection from finite noisy observations without access to the target's exact Fourier coefficients. The forward model was the eight-input algebraic tree used in the preceding representability analysis. Each of its seven internal nodes used the four-coefficient local computation in Eq.~\ref{eq:s_multiaffine_node}, giving 28 fitted coefficients, 14 edges and a root-only output. Leaves received individual named inputs; no global input mixer, additional decoder or bypass was available. This experiment concerns that constrained multilinear model, rather than conductance dynamics or biologically local plasticity.

Tasks were drawn from four equally weighted families: sums of four disjoint pair products, sums of two complementary quartic products, sums of nested products of degrees 2, 4, 6, 8, and a random-interaction control. The control selected four distinct nonconstant monomials uniformly from all degree-2–8 masks, independently of the candidate trees. Each task received a random input permutation where applicable, independent random coefficient signs and magnitudes sampled uniformly between 0.5 and 1.5. Coefficients were normalized to unit squared norm. Under independent Rademacher inputs, each target therefore had mean zero and population variance one. These random coefficients generated new functions; they did not preserve an identical input-gradient second moment across tasks. The earlier exact capacity experiment provides the separate isospectral comparison.

The primary candidate pool retained the same twelve trees: three shapes (balanced, comb and 3+5) under four fixed input assignments. An $\ell_1$-penalized linear estimator (Lasso) received input/output observations and the dictionary of all 255 nonconstant Walsh monomials. It received no family label, planted interaction support, true coefficient, teacher derivative or full-domain truth table. Its prior information was the known eight-variable Rademacher domain and a sparse multilinear dictionary. The intercept was fitted separately and excluded from the interaction score, because a global constant does not constrain subtree compatibility.

Calibration comprised 64, 256 or 1,024 independent input draws, with independent Gaussian label noise of standard deviation zero or 0.5. The primary condition was 256 queries with noise standard deviation 0.5. The first 75\% of observations were available for fitting the Lasso and short-pilot trees; the remaining 25\% formed a held-out calibration subset used to choose among the pilot fits (the calibration gate). Thus the primary estimator used 192 observations, while the shared total calibration-query budget was 256. Let $n_{\rm fit}$ be the number of fitting observations and $\operatorname{sd}(y_{\rm fit})$ their label standard deviation. The Lasso penalty was $c\,\operatorname{sd}(y_{\rm fit})\sqrt{2\log(256)/n_{\rm fit}}$. Fitting used cyclic coordinate descent, tolerance $10^{-8}$ and at most 20,000 iterations. The scale $c$ was selected from 0.15, 0.35, 0.65 using five development seed blocks. The frozen value was 0.15. The same development outcomes fixed the best constant candidate and a lookup rule assigning a candidate to the estimated input-gradient rank needed to explain 99.5\% of its energy. Both reduced to \texttt{balanced\_p0}; bins with fewer than two development tasks used the constant fallback.

For each candidate subtree, we matricized the estimated coefficients across the subtree/complement split, removed the row corresponding to the constant function inside the subtree, and computed the squared singular-value tail beyond its first component. The selector minimized the maximum normalized tail. Scores within $10^{-10}$ were treated as ties; ties were resolved by the sum of overlapping cut tails and then by lexicographic candidate name, using the same tolerance. The maximum is a population-MSE lower bound when computed from the true coefficients. Under noisy estimation it is a selection score, not a certified lower bound on the true task. The summed tails are a tie-breaking heuristic and are not a bound.

Baselines comprised fixed representatives of each shape, the development-best fixed candidate, the development-trained estimated-rank-only rule, uniform random choice averaged exactly over the candidate pool, and a two-sweep fitting pilot. The pilot independently fitted every candidate on the same calibration-fit observations from one fixed initialization and selected the smallest MSE on the calibration gate. Pilot parameters were discarded before final training, so the selected candidate received no additional training from this procedure. A secondary selector applied the previously audited subset dynamic program to the estimated coefficients at the primary calibration condition only. It could select a tree outside the twelve-candidate pool while retaining the same coefficient and edge budget.

All model definitions, estimator choices, baselines, training settings, primary comparisons and effect-size margin were frozen after development. Twenty new seed blocks generated one new function from each family, giving 80 confirmatory tasks. Each task was evaluated under all six calibration conditions. Cryptographic hashes recorded every calibration score, pilot choice, adaptive-tree choice and calibration dataset before any confirmatory candidate completed final training or exposed its final test endpoint. Two implementation errors stopped workers before fitting: lexicographic JSON ordering of numerical node IDs, and a launcher importing a historical module with the same name. The failure logs were retained. Separately hashed amendments corrected only node decoding and module dispatch. The original frozen selector files and all sealed choices remained unchanged.

For each task, final candidate training used a separate fixed stream of 1,024 input draws with Gaussian label noise of standard deviation 0.3. Every candidate used 32 alternating least-squares sweeps and two paired random initializations. Coordinate updates solved the four-parameter local least-squares subproblem, with output-preserving gauge normalization. A separate 512-observation noisy validation stream selected between the two final restarts. No test observation was used for fitting, restart choice or calibration. The short pilot and final fits used separate data and began independently from their declared initializations.

Final evaluation used 4,096 independent test inputs. The primary endpoint was MSE against the clean target, divided by its known population variance of one. MSE against independently noisy test labels and exact MSE over all 256 input patterns were secondary evaluation-only endpoints. Because the input domain is finite, independent streams can contain the same input pattern. They represent independent observations with fresh noise and distinct sample identifiers; they are not partitions of a reused noisy-label cache. The task generator knows the target coefficients to draw observations, but selectors receive only their stated calibration inputs. Complete coefficients and full-domain values are reserved for post-seal target-informed references and evaluation.

The twelve fixed candidates, the estimated adaptive tree and a privileged true-coefficient adaptive reference were each fitted twice per task, giving 2,240 confirmatory fits. The same final candidate outcomes were reused across calibration conditions. Primary regret was the selected candidate's independent-test normalized MSE minus the retrospective minimum among the twelve fixed candidates, each represented by its validation-selected restart. The adaptive selector's regret against this fixed pool can be negative; such values were retained and describe an advantage from searching a larger tree space. The true-coefficient cut selector, true-coefficient adaptive tree and retrospective best trained candidate in the pool were explicitly privileged references, not deployable baselines.

Uncertainty used 20 seed blocks. The four family-specific paired differences were averaged within each seed before bootstrapping 10,000 seed samples. The two primary comparisons were the estimated-cut selector against the development-best fixed candidate and against the two-sweep pilot, at 256 queries/noise 0.5. Descriptive intervals used 95\% coverage. Bonferroni-adjusted two-sided 97.5\% intervals were used for the two primary comparisons. Meaningful superiority required both a positive adjusted lower bound and mean improvement of at least 0.01 normalized MSE. Other calibration sizes, the noise-free conditions, family-specific results and adaptive construction were prespecified secondary analyses. Failed fits and convergence warnings were retained; nonfinite or linear-algebra failures received the predeclared loss of 1e6 and were counted explicitly.

The prespecified task mixture defines the scope of inference. Confirmation tests fresh functions and observations within that mixture; it does not establish transfer to unknown task families, a scalable compiler at larger input counts, or optimal biological dendritic morphology. The fixed Walsh dictionary grows exponentially with the number of inputs, and the tree fitter is an alternating least-squares procedure rather than a synapse-local plasticity rule.

Supplementary Fig.~\ref{fig:supp_morphology_calibration} reports the complete
calibration sensitivity, family-specific effects, primary adjusted contrasts
and implementation cost. Section~\ref{note:morphology_end_to_end} tests
subsequent gradient learning on estimated architectures. The separate study
in Section~\ref{note:morphology_credit_learning} uses oracle-compatible trees
to examine task, assignment and feedback-resolution effects.

\subsection{Direct confirmation from noisy morphology estimates to gradient learning}
\label{note:morphology_end_to_end}

The calibration study evaluated selected trees with an alternating least-squares
fitter. A separate cohort tested whether its estimates also supported
ordinary gradient learning, completing the sequence from noisy calibration
to tree construction and parameter updates (Supplementary Fig.~\ref{fig:supp_morphology_calibration}G,H). Twenty new seed
blocks (930910--930929) were disjoint from the preceding calibration and
credit cohorts. Each generated one fresh matching, quartic, nested-prefix
and random-interaction target using the same prespecified mixture of signs,
coefficient magnitudes and input assignments. Every target had population
variance one; the input-gradient second moment was not held identical across
this mixture.

We retained the 255-monomial Lasso with penalty scale 0.15, the same
intercept treatment and the same subset dynamic program. The selector received 192 noisy observations
from a total calibration budget of 256, with label-noise SD 0.5. The remaining
64 observations were reserved and were not used to select or tune this
follow-up. Neither family labels nor true support or coefficients were
supplied to the estimator. All eighty estimated trees were recorded together with cryptographic
hashes before parameter learning. The fixed balanced tree was the earlier
development-best baseline. A third tree was constructed from true
coefficients after sealing as a privileged structural reference. No
construction weights or other target-derived initialization were used.

Each tree was trained with either exact path credit or unit root broadcast,
giving six conditions per task and 480 fits. Root sensitivity remained one;
broadcast replaced the six other internal sensitivities by one. No oracle
projection coefficients or proximal-zone definition were needed. Each
forward model had seven multi-affine internal units, 28 coefficients,
fourteen edges and root-only output; trees with a leaf directly below the
root were permitted.

Adam used $(\beta_1,\beta_2)=(0.9,0.999)$ and $\epsilon=10^{-8}$. Learning
rates were copied without retuning from the previous credit-development
record: 0.01 for exact credit and 0.003 for broadcast. The comparison therefore tests frozen optimizer recipes at different
learning rates. All six conditions
shared a label-independent Gaussian initialization with SD 0.5, zero biases,
and the same minibatch indices. Initial coefficients and subsequent updates
were clipped to $[-2,2]$; each condition's gradient norm was clipped at ten.
Learning differentiated half-MSE divided by the known target variance,
using 1,024 updates of 64 observations sampled with replacement from a fixed
2,048-observation training dataset. Training label-noise SD was 0.15.
Separate 4,096-observation test data had independent noise of the same SD.
Calibration, training, test, initialization and minibatch streams each used
a separate deterministic random-seed assignment. Repeated Rademacher input patterns were
allowed; observation identities and noise draws were independent across
splits. There was no validation selection, early stopping, new rate search
or continuation from calibration fits.

The primary endpoint was noisy-test NMSE after 1,024 updates, with expected
noise floor 0.0225 for every family. Clean independent-test and exact
256-pattern population errors were secondary. Two primary paired contrasts
compared fixed minus estimated tree error under exact credit, and broadcast
minus exact error on estimated trees. Four family differences were averaged
within each seed before 10,000 whole-seed bootstrap resamples. The two primary
contrasts used Bonferroni-adjusted two-sided 97.5\% intervals and the
prespecified 0.01-NMSE meaningful-improvement margin. Family-specific and
true-coefficient-reference comparisons used descriptive pointwise 95\%
intervals. The primary pooled improvements were 0.45913 NMSE
$[0.38489,0.53227]$ and 0.60750 $[0.56916,0.64610]$, respectively.

Matching, quartic and nested targets approached the noise floor under the
estimated-tree exact-credit recipe, while random interactions retained
substantial error. Broadcast performed slightly better on matching tasks;
exact credit did not dominate every family. The tree constructed from the full target is a privileged structural
reference, not an oracle for the final endpoint: a better interaction score
need not produce a better finite-training outcome. These results test the declared finite-data
construction and learning pipeline, with its specified function class,
optimizer and credit assumptions, rather than a general morphology law or a
biophysical coefficient mechanism.

A pre-training interface mismatch in an oracle helper stopped the initial
array before initialization or fitting. Retained failure records and a
hashed runtime adapter document the repair; scientific settings and sealed
estimates were unchanged. An excluded development smoke (seed 721000)
subsequently exercised all six conditions through the complete training
horizon. Its six fits are excluded from the fresh sample. All 480 fresh
trajectories and final weights were retained. Independent replay reproduced
all eighty Lasso-to-tree choices from sealed observations and all reported
test, clean-test and population errors from saved final weights, with
maximum metric difference zero. Initialization and source-role hashes
matched, and every parameter remained in the declared box. The complete protocol, runtime adapter, excluded implementation check,
selection record and validation results
are supplied under \texttt{source\_data/morphology\_calibration/end\_to\_end}.


\subsection{Learning compatible trees with exact and restricted credit}
\label{note:morphology_credit_learning}

To separate credit assignment from uncertainty in structure estimation, this experiment supplied a tree compatible with the full target. Three development seeds and twenty fresh seeds generated independent signed and permuted matching, quartet and nested-interaction targets. The complete target Fourier tensor supplied a compatible tree through the preceding dynamic program; no target-derived coefficients initialize learning. An independent permutation of leaf labels preserved the exact shape, 28 coefficients and 14 edges. Chance compatible permutations were retained. The primary matching and quartet families retained identical input-sensitivity spectra, while the nested family had the previously stated anisotropic spectrum.

All seven internal units initialized their non-bias coefficients independently from $\mathcal N(0,0.5^2)$, with zero biases. The same initial coefficients, training cache and minibatch draws were used across paired assignments, credit rules and optimizers. Local features were $(1,u_L,u_R,u_Lu_R)$ and the exact path derivative $q_v=\partial u_{\rm root}/\partial u_v$ supplied credit. Exact gradients agreed with central finite differences for every coefficient in all three families. Root sensitivity was one under every rule. The six nonroot internal units received either the exact field, a unit root broadcast, the orthogonal projection onto a single constant profile, or the projection onto two proximal-subtree indicator profiles. A size-matched permutation of the latter profiles served as a control. All projected coefficients were computed from the current exact field and are therefore oracle coefficients; the experiment did not learn a circuit estimator for them.

Training used 2,048 independently sampled Rademacher examples with output noise SD~0.15 and 1,024 minibatch updates of size~64. The gradient corresponded to half-MSE divided by the known clean target variance. Stochastic gradient descent (SGD) and Adam ($\beta_1=0.9,\beta_2=0.999$, $\epsilon=10^{-8}$) each tested learning rates 0.003, 0.01 and 0.03. Update gradients were clipped at norm~10 and parameters projected to $[-2,2]$, which contains all exact constructions. Every clipping/projection count was recorded. These are constrained finite-horizon training results, not estimates of unrestricted optimization optima.

A separate learning rate for each optimizer and rule minimized mean development population NMSE over all families and both assignments. A common rate per optimizer was also selected on development. All choices were recorded with cryptographic hashes before any fresh fit; all three rates remain in the released outcomes. Every SGD rule selected the upper tested rate, so the grid does not establish its best attainable optimizer setting. The independent test set contained 4,096 examples with fresh SD~0.15 output noise. Test NMSE, normalized by the clean target variance, was primary; noiseless population NMSE on all 256 inputs was secondary. Evaluation at steps 0, 1, 16, 64, 256 and 1,024 retained 21,600 checkpoint rows from 3,600 fresh fits. The endpoints at the development-selected rates contained 1,200 fits. No run was excluded. Ten thousand bootstrap samples of twenty whole paired seed blocks were drawn for pointwise 95\% intervals, averaging families within seed for pooled contrasts. Multiple subgroup intervals are descriptive, not separately multiplicity-adjusted discoveries.

Exact credit strongly improves learning of compatible quartet and nested targets under both frozen optimizer recipes. Broadcast is already effective for matching targets and outperforms exact credit there under SGD. Restricted two-subtree projections help in some Adam conditions but do not consistently beat shuffled profiles, and have essentially no pooled advantage over broadcast under SGD. For compatible trees trained with exact credit and Adam, mean test NMSE was 0.0329; permuting the input labels increased it to 0.6080. At the common Adam rate 0.01, the compatible exact/broadcast averages remain 0.0329/0.5973. Supplementary Fig.~\ref{fig:supp_morphology_credit} reports the full family, optimizer, profile, assignment and rate controls. Numerical outcomes and source/selection hashes are in \texttt{source\_data/morphology\_credit/}.

\subsection{Input grouping and restricted credit in directed conductance trees}
\label{sec:conductance_grouping_bridge}

We used a minimal directed shunting model to test input grouping while holding the branching shape and parameter counts fixed. Four terminal compartments fed two proximal compartments and a soma. Every model had seven compartments, six child couplings, four excitatory contacts, six inhibitory contacts and sixteen trainable conductances. Leak conductance was one; excitatory and inhibitory reversal potentials were one and zero. All compartment activations were the identity, and the soma had no direct synaptic inputs or additional decoder. Let $i\in\{0,1,2,3\}$ index terminal compartments and $j\in\{1,2\}$ proximal compartments. Write $T_i$, $P_j$ and $S$ for terminal, proximal and somatic voltages. The pair $\mathcal S_j$ contains the terminals assigned to proximal compartment $j$. Terminal $i$ has synaptic conductances $g_i^{\mathrm E}$ and $g_i^{\mathrm I}$, and $b_{ji}$ couples it to proximal compartment $j$. The maximal proximal shunt conductance is $h_j$, and $a_j$ couples that compartment to the soma. For terminal activities $x_i^{\mathrm E},x_i^{\mathrm I}$ and proximal shunt input $z_j$,
\begin{align}
T_i&=\frac{g_i^{\mathrm E}x_i^{\mathrm E}}
 {1+g_i^{\mathrm E}x_i^{\mathrm E}+g_i^{\mathrm I}x_i^{\mathrm I}},\\
P_j&=\frac{\sum_{i\in\mathcal S_j}b_{ji}T_i}
 {1+\sum_{i\in\mathcal S_j}b_{ji}+h_jz_j},\\
S&=\frac{a_1P_1+a_2P_2}{1+a_1+a_2}.
\end{align}
The positive local input $z_j$ controls shunt conductance; it carries no teaching signal. Conductances were parameterized as $g=\exp(\theta)$, with $\theta\in[-5,5]$. The matched resources were compartment, contact, coupling and parameter counts and the admissible conductance range; total learned conductance was not constrained to be equal. No fresh run reached a parameter bound. The NumPy implementation was checked compartment by compartment against the repository's production \texttt{forward\_branch\_dynamics} function in its identity-activation, zero-numerical-offset configuration. All sixteen log-conductance derivatives and ten input derivatives agreed with autograd to the numerical tolerances reported in the accompanying tests. This is a directed steady-state model, not a reciprocal cable or spiking model.

The three groupings were $(01\mid23)$, $(02\mid13)$ and $(03\mid12)$. A planted teacher used the same conductance equations, with independent Gaussian perturbations of standard deviation $0.25$ in log conductance around $g^{\mathrm E}=2$, $g^{\mathrm I}=0.7$, $h=3$, $b=3$ and $a=4$. Independent student perturbations had standard deviation $0.4$ around these centers and were shared across paired candidate groupings, credit rules and optimizer settings. Student initialization did not use teacher coefficients. Teacher tasks permuted terminal excitatory/inhibitory input pairs together; the proximal inputs retained their identities. Paired task inputs were inversely permuted so that the realized target values were identical across the three tasks. Thus each task had the same input distribution and the same spectrum of the input-gradient second moment. The latter was verified from 2,048 paired probes per seed; its matrix entries could differ by the corresponding coordinate permutation. Compatibility meant that the student's input grouping matched the teacher's grouping. These tasks are planted mechanistic positive controls within one model class, rather than tests of transfer to an unseen class of biological computations.

Each input was $\exp(\sigma Z-\sigma^2/2)$ for independent standard-normal $Z$. We used $\sigma=0.7$ for terminal excitatory and inhibitory inputs and $\sigma=1.2$ for proximal shunt inputs. All activities and conductances were strictly positive, all voltages lay between the reversal potentials, and there was no label noise. Each task used 1,024 training, 1,024 validation, 4,096 test and 256 calibration examples from independent random streams. Within a task, every candidate and credit rule received the same examples and minibatch indices. Training minimized half-MSE divided by the training target variance. Reported NMSE is MSE divided by the target variance of the evaluation split; it is dimensionless and is not a percentage.

Four feedback conditions shared the forward model. Let $V_n$ be the voltage of compartment $n$ and $L$ the training loss. Write $\rho_n=\partial S/\partial V_n$ for its current exact path sensitivity and $\delta_0=\partial L/\partial S$ for the somatic loss derivative. Exact transport supplied $\delta_0\rho_n$. A fixed calibrated broadcast supplied $\delta_0\gamma_n$, where $\gamma_n$ was the mean initial path sensitivity of that student's own forward states on the label-free calibration inputs. The profile was frozen during training. A one-profile projection used the same $\bm\gamma$ and supplied its best Euclidean fit to the current nonsomatic field. The two-subtree condition restricted $\bm\gamma$ to each proximal compartment and its two terminal descendants, giving disjoint columns $\bm A_j$. Its coefficients were
\begin{equation}
c_j=\frac{\bm A_j^{\mathsf T}\bm\rho}
 {\bm A_j^{\mathsf T}\bm A_j}.
\end{equation}
The one-profile projection used the analogous coefficient with $\bm A=\bm\gamma$. These per-trial coefficients had access to the current exact student path field; both projection rules are oracle diagnostics of restricted feedback, not learned coefficient estimators. Soma error remained exact in every condition. Each conductance update multiplied the delivered compartment error by its local conductance eligibility, including the chain derivative of the exponential parameterization. For a child coupling, the local conductance derivative was $(V_{\rm child}-V_{\rm parent})/G_{\rm parent}^{\rm tot}$, where $G_{\rm parent}^{\rm tot}$ is the parent's total open conductance. Credit geometry was evaluated at the corresponding exact-learning model's common checkpoint state; the reported path-capture normalization excluded the soma.

We used 1,000 minibatch updates of size 128, with checkpoints at 0, 10, 100, 300 and 1,000. Stochastic gradient descent (SGD) and Adam were evaluated separately; Adam used $(\beta_1,\beta_2)=(0.9,0.999)$ and $\epsilon=10^{-8}$. One learning rate per optimizer was chosen by averaging final validation NMSE over five development seeds (15300--15304), all task/student groupings and all four credit rules. The original SGD grid was $0.001,0.003,0.01$; because its optimum was at the upper boundary, an additional development-only bracket tested $0.03,0.1,0.3,1$. Every result, including unstable high-rate trajectories, was retained. Adam tested $0.003,0.01,0.03$. The selected rates were $0.3$ for SGD and $0.01$ for Adam. These settings were frozen before the twenty fresh seeds (15400--15419), producing $20\times3\times3\times4\times2=1,440$ fits and 7,200 checkpoint outcomes. No learning rate, horizon or credit rule was selected using fresh training outcomes. The two primary contrasts were the Adam exact-gradient grouping effect and the Adam calibrated-broadcast-minus-exact effect within compatible groupings. Other contrasts are descriptive optimizer and feedback-resolution sensitivities.

Inference used the seed as the independent unit. Each seed's compatible mean averaged its three teacher-task permutations; its incompatible mean additionally averaged the two mismatched student groupings per task. Confidence intervals used 10,000 bootstrap resamples of these paired seed summaries. Intervals are pointwise rather than simultaneous across all reported subgroup comparisons. All sources, rates, states, physical-domain checks and numerical tests are retained under \texttt{source\_data/morphology\_conductance}; execution code is under \texttt{scripts/morphology\_conductance}.

\paragraph{An interaction lower bound for this conductance model.}
Let $R_i$ denote a teacher terminal response, $A_j=(1+\sum_i b_{ji}+h_jz_j)^{-1}$ its proximal attenuation, and $\beta_j=a_j/(1+a_1+a_2)$. The sum in $A_j$ runs over the teacher terminals assigned to proximal compartment $j$. The teacher output is a sum of terms $\beta_j b_{ji}R_iA_j$. Let $p_{\rm teacher}(i)$ and $p_{\rm student}(i)$ denote the proximal parent assigned to terminal $i$ in the teacher and student, respectively, and write $p(i)=p_{\rm teacher}(i)$. For terminal $i$, set $j=p(i)$. Because terminal inputs and proximal shunt inputs are independent, a term whose terminal and teacher shunt belong to different \emph{student} proximal input blocks has the centered interaction
\begin{equation}
H_i=\beta_jb_{ji}(R_i-\mathbb E R_i)(A_j-\mathbb E A_j).
\end{equation}
Every admissible student output is additive across its two proximal input blocks. Each such $H_i$ is orthogonal to all functions additive across those blocks, and the omitted $H_i$ are mutually orthogonal. Therefore, for every choice of the student's conductances,
\begin{equation}
\mathbb E[(S_{\rm teacher}-S_{\rm student})^2]
\geq\sum_{i:\,p_{\rm student}(i)\ne p_{\rm teacher}(i)}
 (\beta_{p(i)}b_{p(i)i})^2
 \operatorname{Var}(R_i)\operatorname{Var}(A_{p(i)}).
\end{equation}
The bound is zero for compatible groupings. It is a lower bound for the larger class of arbitrary functions additive across the student's proximal blocks, not a claim that the constrained rational student attains it. Its validity relies on the specified independent inputs, identity somatic activation and absence of direct somatic synaptic inputs; it does not extend automatically to arbitrary conductance architectures or nonlinear readouts. An equivalent obstruction is the absence of a mixed derivative between a terminal input and a shunt input assigned to the other student subtree.

We specified this diagnostic before the fresh training outcomes and retained the learner and rate-selection procedure unchanged. Population lognormal moments were evaluated with 32-, 64- and 128-node Gauss--Hermite quadrature. The primary bound used 128 nodes; the maximum change from 64 to 128 nodes was below $8\times10^{-17}$ in normalized-MSE units. This is a converged numerical evaluation of an analytic bound, not a formally certified quadrature error. Across the fresh teacher draws, the mean incompatible population-NMSE lower bound was $0.013222$, with a minimum of $0.005367$ across task/candidate pairs. These are population lower bounds, distinct from held-out sample NMSE.


\subsection{Finite-horizon forecasts for the contextual linear learner}
\label{note:finite_horizon_morphology}

To diagnose the original scalar extrapolation independently of structural
capacity, we froze a new finite-horizon predictor before fresh calibration
or training. The candidate pool, learning rate $\eta=0.35$, minibatch size
$B=32$, 256-update horizon and objective cost remained unchanged. Five
original development seeds were diagnostic. Twenty new seeds (9300--9319)
crossed four ranks, two noise levels, two held-out angles, both transfer arms
and both training regimes. Every candidate was trained, yielding 25,600 new
fits and 76,800 checkpoint outcomes at updates 16, 64 and 256. The new
non-pilot scores were hash-sealed before each seed's outcomes. No predictor
hyperparameter was fitted to the new endpoints.
This separates mode-dependent learning dynamics from a scalar initial slope,
as motivated by analyses of deep linear learning and task alignment with
kernel eigenfunctions \citep{lampinen2019generalization,canatar2021spectral}.

To forecast learning, we express the weights in route coordinates. For
orthonormal routes $\Phi$, the learner stays in $W=\Phi V$. For context
$j$ with decoder $\bm a_j$, define the effective route vector
$\bm c_j=\Phi^{\mathsf T}H_T^{\mathsf T}\bm a_j$.
The scalar prediction is $\bm c_j^{\mathsf T}V\bm x$ and the actual update
is minibatch descent in $V$. Let $\bm v=\operatorname{vec}(V)$ stack the route-coordinate weights.
For calibration example $i$ with input $\bm x_i$, label $y_i$ and effective
route vector $\bm c_i$, define
$\bm z_i=\operatorname{vec}(\bm c_i\bm x_i^{\mathsf T})$. With
$\E_i$ denoting an average over the calibration-fit examples, set
$H_{\rm emp}=\E_i[\bm z_i\bm z_i^{\mathsf T}]$ and
$\bm b_{\rm emp}=\E_i[y_i\bm z_i]$. Write
$H_{\rm emp}=U\operatorname{diag}(\lambda)U^{\mathsf T}$, with orthonormal
eigenvectors in $U$ and eigenvalues in $\lambda$. After $T$ updates,
full-batch descent from zero gives
\begin{equation}
 \bm v_T=U\operatorname{diag}\!\left(
 \frac{1-(1-\eta\lambda)^T}{\lambda}\right)U^{\mathsf T}\bm b_{\rm emp},
 \label{eq:s_spectral_forecast}
\end{equation}
The diagonal expression is evaluated separately at each eigenvalue, with
continuous value $\eta T$ when that eigenvalue is zero. The empirical selector
fits the first 512 calibration examples and evaluates on the other 512;
Eq.~\ref{eq:s_spectral_forecast} is exact for that full-batch surrogate,
not for the original cached-data SGD process.

For fresh inputs $\bm x\sim\mathcal N(0,I_d)$, where $d$ is the input
dimension, a smaller recurrence describes stochastic gradient descent (SGD).
Let $p_j$ be the context probabilities and $\bm t_j\in\mathbb R^d$ the
target coefficient vector in context $j$. Suppose
$y=\bm t_j^{\mathsf T}\bm x+\epsilon$, with independent zero-mean noise
of variance $\sigma_j^2$. The exact recurrence closes on the mean weights
$M=\E[V]\in\mathbb R^{K\times d}$ and the route-space variance matrix
$\Gamma=\E[(V-M)(V-M)^{\mathsf T}]$. The latter sums weight covariance
over the input coordinates; it is not the full vectorized parameter
covariance. Define
\begin{align}
 C&=\sum_jp_j\bm c_j\bm c_j^{\mathsf T},&
 D&=\sum_jp_j\bm c_j\bm t_j^{\mathsf T},&
 \bar G&=CM-D,\\
 R_j&=\|M^{\mathsf T}\bm c_j-\bm t_j\|^2+
                    \bm c_j^{\mathsf T}\Gamma\bm c_j,&
 A_\eta&=I-\eta C,\\
 J&=\sum_jp_j\bm c_j\bm c_j^{\mathsf T}
                  [(d+2)R_j+d\sigma_j^2],&
 K_G&=C\Gamma C+\bar G\bar G^{\mathsf T}.
\end{align}
Here $\bar G$ is the mean gradient, $R_j$ is the expected squared residual
excluding observation noise, $J$ is the second moment of the single-example
gradient Gram matrix, and $K_G$ is the corresponding second moment of the
conditional-mean gradient. For a fixed residual coefficient vector
$\bm e$, Gaussian fourth moments yield
$\E[(\bm e^{\mathsf T}\bm x-\epsilon)^2\|\bm x\|^2]
=(d+2)\|\bm e\|^2+d\sigma_j^2$. Conditioning first on the current
weights, the Gram of a minibatch-mean gradient is its conditional-mean Gram
plus $1/B$ times its conditional variance. Total covariance therefore gives
\begin{align}
 M_{t+1}&=M_t-\eta\bar G_t,\\
 \Gamma_{t+1}&=A_\eta\Gamma_t A_\eta^{\mathsf T}
                      +\frac{\eta^2}{B}(J_t-K_{G,t}),
 \label{eq:s_fresh_sgd_moments}\\
 \E[\mathcal L_T]&=\tfrac12\sum_jp_j[R_j(T)+\sigma_j^2].
\end{align}
Thus the forecast updates gradient noise as the residual and weight
variance change during learning. Gaussian input fourth moments and fresh independent
minibatches are essential. Noise need only have the stated first two moments.
This closure is not asserted for nonlinear targets, anisotropic inputs,
adaptive optimizers or correlated samples.

The primary plug-in predictor estimates context frequencies, conditional
coefficient vectors and residual noise variances using 1,024 independent
calibration examples. Coefficients use ordinary least squares without an
intercept, and noise variances use residual degrees of freedom. Its recurrence
is exact for the fitted Gaussian population and is a plug-in approximation
for the generating population. Population-oracle versions instead receive
true teacher coefficients and noise moments and are labeled privileged.
The 16-step pilot performs real candidate updates and evaluates on
calibration observations; its computation is included. An observed-context-
count baseline chooses the smallest permitted $K$ covering the number of
distinct observed context vectors, followed by the cheapest tree. This baseline and a conditioning baseline restricted to the same observed
count were added after inspecting development results, then frozen
separately before the new outcomes. Context count equals generating
rank in this constructed family; it is not a general estimator of credit rank.

We ran both fresh-example SGD and minibatches resampled from a fixed
2,048-example cache, retaining a common 4,096-example test set. Reusing a
cache makes later weights dependent on those examples and violates the fresh
recurrence's independence assumption. Population-oracle prediction minus
exact population loss of the realized weights quantifies that discrepancy.
All task conditions are averaged within each independent seed before 10,000
whole-seed bootstrap draws. Intervals are pointwise; the many secondary
selector, rank and regime comparisons are not interpreted as separately
confirmed discoveries. No new multiplicity-adjusted significance claim is
made from these interval plots. Undefined correlations in numerically
constant cases are retained with explicit defined-task counts rather than
used to exclude outcomes.

Supplementary Fig.~\ref{fig:supp_finite_horizon_selection} shows endpoint
regret, strong baselines and implementation time;
Supplementary Fig.~\ref{fig:supp_finite_horizon_prediction} shows prediction
accuracy and cache sensitivity. The primary forecast reaches mean fixed-cache
regret $0.000064$ and $0.000042$ in the feedback and joint arms, while the
cheap observed-count baseline reaches $0.000134$ and $0.000197$. Thus the
large improvement over the original scalar score does not establish a large
advantage for fine morphology selection. Focused tests compare singular
spectral formulas with direct descent, stochastic moments with exact Gaussian
quadrature, and enforce that predictors cannot request training or test
observations. Checks of source hashes and pre-outcome selection records passed. Numerical data,
protocols, complete outcomes and exported source hashes are provided in
\texttt{source\_data/morphology\_structure/} and
\texttt{source\_data/morphology\_finite\_horizon/}.
